\section{Introduction}
\label{sec:introduction}

Every organisation runs on numbers it cannot fully explain. A revenue figure on a dashboard is the
end of a chain: a loader lands rows from a source system, dbt models clean and join them, an
Airflow task copies the result into a reporting table, and a semantic model in a business
intelligence tool aggregates it. Each link is code, written by different people in different
tools, and the meaning of the final number (which rows count, in which currency, net of what)
lives partly in that code and partly in people's heads. When the number is questioned, someone
reads the chain by hand.

The cost of that is measured, if imperfectly. Gartner puts the average cost of poor data quality
at USD 12.9 million a year per organisation, from research published in 2020 and based on
self-estimated costs~\citep{gartner2020dq}. In a 2024 vendor survey of over 550 respondents, 67\%
said they do not completely trust the data they use for decisions~\citep{precisely2024}. In a 2023
vendor survey of 200 data professionals, a data incident took 15 hours on average to resolve, up
166\% on the year before~\citep{montecarlo2023}. dbt Labs' 2026 vendor survey of 363 data
professionals found 71\% concerned about incorrect data reaching stakeholders and 41\% naming
ambiguous data ownership as an obstacle~\citep{dbt2026state}. Software developers, in a 2018 field
study, spent 58\% of their time on program comprehension~\citep{xia2018comprehension}, and a data
stack is a program spread across many tools. Regulators ask for the chain explicitly: in 2023, ten
years after the Basel Committee's principles for risk data aggregation, only 2 of 31 global
systemically important banks were fully compliant with all of them~\citep{bcbs2023}, and Article 10
of the EU Artificial Intelligence Act of 2024 requires the data of high-risk systems to be subject
to data governance covering its origin and preparation~\citep{euaiact}.

AI agents raise the stakes. An agent that writes SQL against a warehouse it does not understand
produces confident, plausible and wrong numbers. On the enterprise questions of Sequeda et al.\
(2023), GPT-4 answered 16.7\% correctly from the SQL schema alone and 54.2\% when the same
questions were posed over a knowledge graph of the business~\citep{sequeda2023benchmark}. On
Spider 2.0 (2025), which asks for real enterprise workflows, the o1-preview agent solves 21.3\% of
tasks, against 91.2\% on Spider 1.0 and 73.0\% on BIRD for GPT-4-based
methods~\citep{lei2025spider2}; on BIRD itself (2023), GPT-4 reached 54.89\% execution accuracy
against 92.96\% for people~\citep{li2023bird}. A 2026 vendor study reports 64.5\% accuracy for
text-to-SQL on a raw normalised schema, 84.1\% to 90.0\% for text-to-SQL on the modelled project,
and 98.2\% to 100\% through a semantic layer over it~\citep{dbt2026semantic}. The lesson of these results is consistent:
agents are only as good as the facts about the data they are grounded in. Those facts are
lineage, formulas, definitions and relationships, and today they are mostly asserted by tools
that never check them.

This paper describes Konstruct, an MIT-licensed system that makes those facts explicit and
checkable. It reads a repository and its warehouse locally, in one command, with no server, no
account and no telemetry, and writes a graph that people and agents query from a command line,
a Python API, an MCP server, a GitHub Action and a single offline HTML page. Its contributions
are:

\begin{enumerate}
  \item \textbf{One graph across tools.} A typed graph model (\cref{sec:graph}) that joins
  column-level lineage from dbt, standalone SQL, Airflow and Dagster tasks, PySpark jobs,
  warehouse views, Looker and Power BI semantic models and the tables Fivetran and Airbyte land,
  through DuckDB, Postgres, Snowflake, BigQuery, Databricks, Trino and Iceberg, without importing
  or running any user code (\cref{sec:chains}).
  \item \textbf{Formulas, not just arrows.} Canonical column expressions over physical columns
  and their composition through every model upstream, so a column reads as a formula over source
  columns, with a text composition that is exact where the SQL printer is operand-insensitive and
  falls back to trees where it is not (\cref{sec:sql,sec:compose}).
  \item \textbf{Lineage measured against the data.} A verification procedure by perturbation
  that confirms claimed edges and exposes missing ones with concrete witnesses, with a soundness
  argument and an honest account of what it cannot see (\cref{sec:verify}).
  \item \textbf{Meaning derived, not declared.} An ontology of entities, measures and
  relationships derived from keys, tests and joins, bound to approved business definitions, with
  ontology packs for common source schemas and round trips of Apache Ossie semantic models
  that return what was authored (\cref{sec:ontology,sec:ossie,sec:packs}).
  \item \textbf{Grounding for agents.} An MCP interface whose answers carry lineage, formulas and
  verification state, and a gate that a change must pass before it is proposed
  (\cref{sec:agents}).
  \item \textbf{A paper regenerated from the code.} Every measurement of Konstruct here is produced by
  \cmd{paper/generate.py} from the code at the commit that built this PDF, and a test fails when
  the paper types a number outside a cited sentence or leaves out a command, tool or format
  (\cref{sec:evaluation}).
\end{enumerate}
